% TOP_PROBLEM: {"id": 1436, "title": "Line Graphs of Non-Word-Representable Graphs", "category_id": 3, "category": {"id": 3, "name": "graph_theory", "display_name": "Graph Theory", "description": "Problems involving graphs, networks, and their properties.", "slug": "graph-theory", "order_index": 3, "created_at": "2026-07-31T15:26:25.671Z"}, "status": "open"}
% FIRST_POSTED: null
% ENTRY_KIND: statement_only
% Imported from: batch07.tex; UnsolvedMath ID 1436
\documentclass[11pt]{article}
\usepackage[margin=25mm]{geometry}
\usepackage{fontspec}
\setmainfont{Latin Modern Roman}
\usepackage{amsmath,amssymb,mathtools}
\usepackage{unicode-math}
\setmathfont{Latin Modern Math}
\usepackage{xurl,hyperref}
\hypersetup{colorlinks=true,urlcolor=blue}
\setcounter{secnumdepth}{0}
\setlength{\parindent}{0pt}
\setlength{\parskip}{5pt}
\setlength{\emergencystretch}{3em}
\newcommand{\sref}[2]{\hyperlink{source:#1:#2}{[S#2]}}
\newcommand{\cataloguescope}[1]{\paragraph{Recorded target.}#1}
\begin{document}
% UnsolvedMath ID 1436; source catalogue code GRAPH-049
\setcounter{section}{1435}
\section{Line Graphs of Non-Word-Representable Graphs}\label{1436}
{\small\sffamily UnsolvedMath ID 1436 · Graph Theory}
\subsection{Definitions and mathematical statement}

\paragraph{Shared notation.}$\mathbb N=\{1,2,\ldots\}$, and $\mathbb Z,\mathbb Q,\mathbb R,\mathbb C$ denote the integers, rationals, reals and complex numbers. Any other conventions are specified locally.


Let $G=(V,E)$ be a finite simple undirected graph. A word over $V$ is a finite sequence of vertices in which every vertex occurs. For distinct $x,y\in V$, let $w|_{\{x,y\}}$ be the word obtained by deleting every letter other than $x,y$. They alternate when this restriction has no equal consecutive letters, that is, it has the form $xyxy\cdots$ or $yxyx\cdots$, of either even or odd length. The word $w$ represents $G$ when, for every distinct $x,y$,
\[
 xy\in E\quad\Longleftrightarrow\quad x,y\text{ alternate in }w.
\]
Write $\mathrm{WR}(G)$ when such a word exists.

The line graph $L(G)$ has one vertex for each edge of $G$; two different vertices of $L(G)$ are adjacent exactly when the corresponding edges of $G$ share an endpoint.
\paragraph{Statement recorded in the supplied source.}
Is the implication
\[
 \neg\mathrm{WR}(G)\quad\Longrightarrow\quad\neg\mathrm{WR}(L(G))
\]
valid for every finite simple graph $G$? Equivalently, can a non-word-representable graph have a word-representable line graph? \sref{1436}{2}
\subsection{Short English statement}
Does taking a line graph always preserve failure of word representation, or is there a graph that fails to be representable while its line graph is representable?
\subsection{Sources}
\begin{description}
\item[\hypertarget{source:1436:1}{S1}] ULAM.AI and the UnsolvedMath Contributors, UnsolvedMath: A Curated Collection of Open Mathematics Problems, record 1436 (GRAPH-049). CC BY 4.0. \url{https://huggingface.co/datasets/ulamai/UnsolvedMath}.
\item[\hypertarget{source:1436:2}{S2}] Khyodeno Mozhui, Tithi Dwary and K. V. Krishna, Line Graphs of Non-Word-Representable Graphs Are Not Always Non-Word-Representable, arXiv:2509.03339v1 (2025), Sections 1 and 3, Theorem 3.1. \url{https://arxiv.org/pdf/2509.03339v1}.
\end{description}
\label{end:1436}
\end{document}
